一张上千列的数据表，未必真有上千个彼此独立的变化方向。多个传感器可能记录同一个物理状态，多项消费指标可能反映少数几种偏好，相邻像素之间的相关强到几乎冗余。这一单元利用的第一性事实就是这条：观测坐标多，不等于生成数据的自由度多。

但请从一开始就避开一个说法：降维不是``把列数减少''。挑几列出来只是把坐标轴限制在原始变量上，而降维允许换一整套坐标；更要紧的是，任何降维方法真正要求你签字的是另一件事——\textbf{哪些方向上的信息可以丢}。这个单元的七篇文章，可以整个读成对这个问题的七种不同回答。

\begin{center}
\begin{tikzpicture}[>=stealth,scale=1.0]
\node[font=\small] at (5.6,2.55) {同一个问题：允许丢掉哪些方向的信息？};
\draw[black!55] (0.6,2.25) -- (10.6,2.25);
\foreach \x in {1.4,4.1,6.9,9.6} \draw[black!55] (\x,2.25) -- (\x,1.95);
% 四个盒子
\draw[black!70] (0.1,0.55) rectangle (2.7,1.95);
\node[font=\scriptsize,align=center] at (1.4,1.6) {丢方差小的方向};
\node[font=\scriptsize,align=center] at (1.4,1.0) {PCA、截断 SVD\\矩阵补全};
\draw[black!70] (2.9,0.55) rectangle (5.3,1.95);
\node[font=\scriptsize,align=center] at (4.1,1.6) {丢全局距离};
\node[font=\scriptsize,align=center] at (4.1,1.0) {Isomap、LLE\\t-SNE、UMAP};
\draw[black!70] (5.6,0.55) rectangle (8.2,1.95);
\node[font=\scriptsize,align=center] at (6.9,1.6) {几乎不挑方向};
\node[font=\scriptsize,align=center] at (6.9,1.0) {随机投影\\（只保点对距离）};
\draw[black!70] (8.4,0.55) rectangle (10.8,1.95);
\node[font=\scriptsize,align=center] at (9.6,1.6) {丢特有噪声};
\node[font=\scriptsize,align=center] at (9.6,1.0) {因子分析、PPCA\\ICA、CCA};
% 底部说明
\node[font=\scriptsize,black!65] at (1.4,0.15) {全局线性};
\node[font=\scriptsize,black!65] at (4.1,0.15) {只信局部线性};
\node[font=\scriptsize,black!65] at (6.9,0.15) {与数据无关};
\node[font=\scriptsize,black!65] at (9.6,0.15) {写成生成模型};
\end{tikzpicture}
\end{center}

\emph{四类方法的分岔点不在算法，而在各自允许丢掉什么。}

顺着这张图，几条贯穿全章的判断值得先摆出来。中心化不是可选的预处理：忘了它，第一主成分会指向数据的重心，症状是解释方差比高得离谱而所有样本得分几乎相同。高维下样本协方差的谱会系统性失真，碎石图上的肘部可能纯由噪声造出，判断信号该拿 Marchenko--Pastur 上界作基准。截断 SVD 的最优性来自 Eckart--Young，而它只在无约束、酉不变范数下成立，加上非负、稀疏、量化或观测掩码之后，截断连候选都算不上。非线性方法保的是邻域不是距离，所以 t-SNE 图上的簇间距离和簇大小都不可解释，超参数一改结果就变。最后，降维一旦接进预测流程，问题本身已经换了，投影必须和模型一起进入交叉验证。

\subsection{七篇的展开顺序}

前三篇共用``低秩''这个词，任务却不同：PCA 为样本变化找一套坐标系，截断 SVD 近似一个已知矩阵，矩阵补全要从部分条目推断整个未知矩阵。把这三件事混作一件，最常见的后果是把一个代数结论当成统计保证。第四篇保留``自由度少''的信念，只放弃``全局线性''，于是得到展开弯曲结构的能力和一整套新的脆弱性。第五篇把方向反过来看，说明升维与降维是分工而非矛盾，并把随机投影这条与数据无关的路线放进来。最后两篇回到生成与可辨认性：同样是线性潜变量模型，采用各向同性噪声、坐标特异噪声、非高斯独立性还是跨视图相关，得到的方向完全不同。

主干路线是 PCA \(\rightarrow\) SVD \(\rightarrow\) 矩阵补全。若关心可视化和局部几何，接着读流形那篇；若关心潜变量的现实含义，跳到因子分析与 ICA、CCA，并且把可辨认性放在解释之前。第一遍读完，能说清``保留方差''、``保留距离''、``最小重构误差''、``恢复未观测条目''为什么是四个不同的目标就够了，谱分解的证明可以留到第二遍。

\subsection{本章篇目}

以下按概念依赖排列；若从具体困惑进入，也可以先打开最接近的一篇，再沿篇内链接回补。

\begin{itemize}
\tightlist
\item \hyperref[sec:13-Machine-Learning/04-Dimensionality-Reduction/01]{``PCA 从最大方差到最小重构误差''}；
\item \hyperref[sec:13-Machine-Learning/04-Dimensionality-Reduction/02]{``SVD 给出最佳低秩近似''}；
\item \hyperref[sec:13-Machine-Learning/04-Dimensionality-Reduction/03]{``矩阵补全从缺失条目到低秩正则''}；
\item \hyperref[sec:13-Machine-Learning/04-Dimensionality-Reduction/04]{``流形假设与非线性降维''}；
\item \hyperref[sec:13-Machine-Learning/04-Dimensionality-Reduction/05]{``降维与升维的对偶：在合适的维度做合适的事''}；
\item \hyperref[sec:13-Machine-Learning/04-Dimensionality-Reduction/06]{``因子分析与概率 PCA：把低秩变成生成模型''}；
\item \hyperref[sec:13-Machine-Learning/04-Dimensionality-Reduction/07]{``ICA 与 CCA：选择不同的可辨认结构''}。
\end{itemize}

七篇之外还有一条不在目录里的线索：每一篇的最后都在划边界，说明该方法在什么条件下不再成立。把这些边界连起来读，比记住任何一个公式都更接近实际工作中要做的判断。
